% =====================================================================
%  Two-screen scintillation formalism and screen-distance constraints
%  Methods / theory section for the CHIME+DSA-110 FRB scattering paper
%
%  Every equation below has been checked against the primary literature
%  (Cordes & Rickett 1998; Masui et al. 2015; Cordes & Chatterjee 2019;
%  Main et al. 2022; Sammons et al. 2023; Pradeep et al. 2025) and, where
%  geometric, against an independent numerical ray-trace. Constants that a
%  first-pass draft got wrong (see notes) are given here in their correct,
%  sourced form.
% =====================================================================

\section{Scintillation formalism and two-screen geometry}
\label{sec:twoscreen}

\subsection{ACF-based diffractive-scintillation estimator}
\label{sec:acf}

For each burst we form the on-pulse spectrum $I(\nu)$ by averaging the
dynamic spectrum over the on-pulse samples, and estimate the noise level
from an off-pulse region. The one-dimensional frequency autocorrelation
function (ACF) of the mean-subtracted spectrum is
%
\begin{equation}
\mathrm{ACF}(\delta\nu)
  = \frac{\big\langle\, \Delta I(\nu)\,\Delta I(\nu+\delta\nu)\,\big\rangle_\nu}
         {\big\langle\, \Delta I(\nu)^2\,\big\rangle_\nu},
\qquad
\Delta I(\nu) \equiv I(\nu) - \langle I\rangle,
\label{eq:acf}
\end{equation}
%
where $\langle\cdot\rangle_\nu$ denotes an average over frequency channels.
The zero-lag point ($\delta\nu=0$) is \emph{excluded} from all fits: radiometer
noise is uncorrelated between channels, so it contributes a delta-function
spike at $\delta\nu=0$ that biases the peak upward by the noise variance.
The scintillation modulation index $m\equiv\sigma_I/\langle I\rangle$ is
recovered as the noise-corrected zero-lag amplitude, i.e.\ the value of
the fitted model extrapolated to $\delta\nu\to0$,
%
\begin{equation}
\mathrm{ACF}(0^{+}) = m^{2}.
\label{eq:m2}
\end{equation}
%
A point source in the strong-scattering regime gives $m=1$; a source whose
apparent angular size partially resolves the scattering disk gives $m<1$
(partial quenching), the property we exploit in
Section~\ref{sec:distance}.

We model each diffractive component with a Lorentzian, the functional form
derived by \citet{Gwinn1998} for a single thin screen,
%
\begin{equation}
\mathrm{ACF}(\delta\nu)
  = \frac{m^{2}}{1 + (\delta\nu/\nu_{s})^{2}},
\label{eq:lorentzian}
\end{equation}
%
where the decorrelation (scintillation) bandwidth $\nu_{s}$ is the
half-width at half-maximum. The scintillation bandwidth and the
pulse-broadening time $\tau_{s}$ produced by the \emph{same} screen are
conjugate through the Fourier uncertainty relation
%
\begin{equation}
2\pi\,\nu_{s}\,\tau_{s} = C_{1},
\label{eq:c1}
\end{equation}
%
with $C_{1}$ a geometry- and spectrum-dependent constant of order unity.
As summarised by \citet{Pradeep2025}, $C_{1}$ ranges over $0.5$--$2$
across the plausible screen-geometry and turbulence-spectrum combinations
\citep{LambertRickett1999}, being of order unity for a thin screen. We do
not adopt a single point value: we take $C_{1}=1$ as the fiducial thin-screen
value and propagate the full $0.5$--$2$ range as a systematic on any
$\tau_{s}\leftrightarrow\nu_{s}$ conversion.
%
% NOTE (not for manuscript): the only C1 statement we can source directly is
% the 0.5-2 range (Lambert & Rickett 1999, as quoted in the fetched
% Pradeep+2025 text). Specific point values (e.g. 1.16 Kolmogorov, 0.74
% uniform medium) were NOT verifiable against retrieved primary text and are
% deliberately omitted; confirm against Cordes & Rickett 1998 Table before
% quoting any single value.

\paragraph{Finite-sample statistics.}
Across an observing band $B$ the number of independent scintles is
%
\begin{equation}
N_{\rm scint} \simeq \eta\,\frac{B}{\nu_{s}},
\qquad \eta = \mathcal{O}(1),
\label{eq:nscint}
\end{equation}
%
with a filling factor $\eta$ that accounts for scintle overlap. The finite
number of scintles sets an irreducible ``self-noise'' floor on the
estimator: the fractional uncertainties scale as
%
\begin{equation}
\frac{\sigma_{\nu_{s}}}{\nu_{s}} \sim \frac{1}{\sqrt{N_{\rm scint}}},
\qquad
\frac{\sigma_{m}}{m} \sim \frac{1}{\sqrt{N_{\rm scint}}}.
\label{eq:selfnoise}
\end{equation}
%
This term dominates the error budget when $\nu_{s}$ is a sizeable fraction
of $B$ (few scintles across the band), and is added in quadrature with the
radiometer-noise contribution when quoting uncertainties on $\nu_{s}$ and
$m$.


\subsection{Two-screen geometry}
\label{sec:geometry}

We place the observer at $O$ and the source at angular-diameter distance
$D$. Two thin screens lie along the sightline: a near (Milky Way) screen at
distance $d_{\rm MW}$ from $O$, and a far (host-galaxy) screen at distance
$d_{h}$, with $d_{\rm MW}<d_{h}<D$. Following the standard thin-screen
result, a screen at distance $d$ imposes a scattering time
%
\begin{equation}
\tau = \frac{\theta_{\rm scat}^{2}\,d_{\rm eff}}{2c},
\qquad
d_{\rm eff} = \frac{d\,(D-d)}{D},
\label{eq:tau}
\end{equation}
%
so the (rms) scattering angle is
$\theta_{\rm scat}=\sqrt{2c\tau/d_{\rm eff}}$, and the diffractive
field-coherence scale at the screen is
$s_{d}=\lambda/(2\pi\theta_{\rm scat})$. (Cosmological factors $(1+z)$ enter
the far-screen terms; we carry them explicitly in
Eq.~\ref{eq:distance} and suppress them here for clarity.)

\paragraph{Apparent size of the host image.}
The quantity that governs whether one screen resolves the other is the
angular size of the host screen's scattering disk \emph{projected to the
near screen}. A first-principles ray-trace (a point source at $D$, a thin
deflector at $d_{h}$, an observing plane at $d_{\rm MW}$) gives the
arrival-direction spread at the near screen
%
\begin{equation}
\theta_{\rm img}
  = \theta_{{\rm scat},h}\,\frac{D-d_{h}}{D-d_{\rm MW}} .
\label{eq:thimg}
\end{equation}
%
% NOTE (not for manuscript): the correct lever-arm denominator is
% (D - d_MW), the far-screen-to-near-screen path measured from the SOURCE
% side. A draft wrote (d_h - d_MW) (the screen-screen separation); that
% fails the d_MW->0 limit and was verified wrong by ray-trace at every d_MW.
%
In the limit $d_{\rm MW}\to0$ (near screen at the observer) this reduces to
the familiar observed image size
$\theta_{h}\simeq\theta_{{\rm scat},h}(D-d_{h})/D
=(2c\tau_{h}\,d_{ls,h})^{1/2}/D$ \citep{Masui2015,Main2022}, with
$d_{ls,h}=D-d_{h}$ the source--host-screen distance.

\paragraph{Resolution power.}
The near screen acts as an interferometer of effective aperture
$L_{\rm MW}$ and angular resolution $\theta_{\rm res}=\lambda/L_{\rm MW}$.
\citet{Pradeep2025} parameterise the degree to which one screen resolves
the other by the \emph{resolution power}
%
\begin{equation}
\mathrm{RP} \equiv \frac{\Theta_{\rm size}}{\theta_{\rm res}}
  = \frac{L_{\rm MW}\,L_{\rm host}}{\lambda\,D_{\rm MW,host}},
\label{eq:rp}
\end{equation}
%
where $L_{\rm MW},L_{\rm host}$ are the transverse scattering-disk sizes at
the two screens and $D_{\rm MW,host}$ the (angular-diameter) distance
between them. $\mathrm{RP}\ll1$: screens do not resolve each other (both
scintillation scales visible); $\mathrm{RP}\simeq1$: marginal resolution;
$\mathrm{RP}\gg1$: the near screen fully resolves the far screen and the
associated broad-scale scintillation is quenched.


\subsection{Modulation-index parameterization for two screens}
\label{sec:modindex}

With two screens the total spectral ACF is \emph{not} simply the sum of the
two Lorentzians. \citet{Pradeep2025} show it carries an additional
\emph{product} term,
%
\begin{equation}
\mathrm{ACF}_{\rm tot}(\delta\nu)
  = \underbrace{m_{\rm MW}^{2}\,\rho_{\rm MW}(\delta\nu)
  + m_{h}^{2}\,\rho_{h}(\delta\nu)}_{\text{sum of screens}}
  + \underbrace{m_{\rm MW}^{2}\,m_{h}^{2}\,
      \rho_{\rm MW}(\delta\nu)\,\rho_{h}(\delta\nu)}_{\text{product term}},
\label{eq:acf2}
\end{equation}
%
where $\rho_{s}(\delta\nu)=[1+(\delta\nu/\nu_{s,s})^{2}]^{-1}$ is the
normalised Lorentzian of screen $s$. In the unresolved regime, with each
screen fully modulated ($m_{\rm MW}=m_{h}=1$), the noise-free zero-lag value
is $\mathrm{ACF}_{\rm tot}(0^{+})=1+1+1=3$, so the combined modulation index
rises to
%
\begin{equation}
m_{\rm tot} = \sqrt{3}\qquad(\text{unresolved, fully modulated}).
\label{eq:sqrt3}
\end{equation}
%
As the near screen begins to resolve the far screen, the broad-scale
component is progressively (not abruptly) quenched: $m_{\rm MW}$ declines
toward zero with increasing RP, the product term collapses, and
$m_{\rm tot}$ falls back toward the single-screen value. The measured
$\{m_{\rm MW},m_{h},\nu_{s,\rm MW},\nu_{s,h}\}$ therefore diagnose the
resolution regime directly, and this is the lever that converts a
scintillation \emph{detection} into a screen-distance constraint.


\subsection{From an upper limit to a distance measurement}
\label{sec:distance}

The physical argument \citep{Masui2015,CordesChatterjee2019} is: the near
(MW) screen can only imprint scintillation if the host image it looks at is
\emph{unresolved}, $\theta_{\rm img}\lesssim\theta_{\rm res,MW}$
(equivalently $\mathrm{RP}\lesssim1$). Observing MW scintillation therefore
bounds the host scattering angle from above, and — because
$\theta_{{\rm scat},h}$ is fixed by the measured host $\tau_{h}$ through
Eq.~(\ref{eq:tau}) — bounds the product of screen distances. In the modern,
redshift-consistent form of \citet{Pradeep2025} (their Eq.~7.6),
%
\begin{equation}
D_{h,\rm FRB}\,D_{\rm MW}
  \;\lesssim\;
  \frac{(1+z_{\rm FRB})\,D_{\rm FRB}^{2}}{8\pi\,\nu^{2}}\,
  \frac{\nu_{s,\rm MW}}{m_{\rm MW}\,\tau_{s,h}},
\label{eq:distance}
\end{equation}
%
where $D_{h,\rm FRB}$ is the (angular-diameter) source--host-screen
distance, $D_{\rm MW}$ the observer--MW-screen distance, $\nu$ the observing
frequency, $\nu_{s,\rm MW}$ the MW decorrelation bandwidth, $m_{\rm MW}$ the
MW modulation index, and $\tau_{s,h}$ the host scattering time. With
$D_{\rm FRB}$ and $z_{\rm FRB}$ fixed by the host redshift and $D_{\rm MW}$
estimated from a Galactic electron-density model (NE2001/YMW16), this is an
\emph{upper limit} on $D_{h,\rm FRB}$ — the distance from the source to its
host screen. Order-of-magnitude it reproduces the
$L_{x}L_{g}\lesssim\text{few}\ \mathrm{kpc}^{2}$ constraints obtained for
CRAFT FRBs \citep{Sammons2023}.

The inequality becomes an \emph{equality} in exactly the marginal case.
Full quenching ($\mathrm{RP}\gg1$) erases the MW scintillation, giving no
measurement; no quenching ($\mathrm{RP}\ll1$) gives only the bound above.
But if the MW scintillation is detected \emph{with reduced modulation}
($0<m_{\rm MW}<1$, i.e.\ $\theta_{\rm img}\simeq\theta_{\rm res,MW}$ and
$\mathrm{RP}\simeq1$), the host image is \emph{marginally} resolved, the
resolution condition is saturated, and Eq.~(\ref{eq:distance}) closes to
%
\begin{equation}
D_{h,\rm FRB}\,D_{\rm MW}
  \;\simeq\;
  \frac{(1+z_{\rm FRB})\,D_{\rm FRB}^{2}}{8\pi\,\nu^{2}}\,
  \frac{\nu_{s,\rm MW}}{m_{\rm MW}\,\tau_{s,h}},
\label{eq:distance_eq}
\end{equation}
%
now a \emph{measurement} of the source--host-screen distance, with a
precision set by how well $D_{\rm MW}$ is known
\citep[cf.][their locating-the-host-screen result]{Pradeep2025}. A measured
$m_{\rm MW}$ that is significantly below unity while the MW scintillation
remains clearly detected is thus the observational signature that promotes
the limit to an equality.

\paragraph{Caveats.} Eqs.~(\ref{eq:distance})--(\ref{eq:distance_eq}) assume
both screens are two-dimensional (isotropic) in the plane of the sky; for
one-dimensional (elongated) screens the relative position angle enters and
the scalar relation breaks down \citep{Pradeep2025}. The quenching of the
broad-scale component is gradual, so ``marginal resolution'' is a regime,
not a knife-edge; we therefore report Eq.~(\ref{eq:distance_eq}) with an
explicit RP-dependent systematic rather than as a point estimate.
